\subsection{What the statute must carry}
\label{sec:11.4.5}
For the statute to start over rather than rename, it has to carry seven terms that neither abolition bill yet contains in full. The terms are mine, not Amnesty's. Where Amnesty asks for something a term carries, the term says so.

\emph{First, the schedule.} Table C.2's fourteen entries are enacted as limitations on all federal funds, binding every agency, contractor and grantee, in the form the Melt ICE Act already uses for detention. Written that way, Border Patrol, the Bureau of Prisons, a FEMA grant and a payment to a foreign government's prison are covered because the act is. The biometric entry needs no new drafting. The moratorium's bar on every federal entity is already in the form the schedule takes \autocite{Jayapal2026Moratorium}.

\emph{Second, succession and separation.} Any office that receives a transferred function inherits the limitations with it. Following section 471(b) of the 2002 Act, no reorganization may recombine in one office the power to arrest, the power to detain, custody of the data used to select people, and adjudication of status. The adjudicator sits outside the department that prosecutes. At present the Board of Immigration Appeals sits in the Justice Department, its members act as the Attorney General's delegates, and the Attorney General may review and replace their decisions \autocite{cfr1003_1}.

\emph{Third, counsel, screening and the stay.} Five provisions make the operator wait for the answer:

\begin{itemize}
\item Anyone facing detention or removal has counsel at government expense from the first encounter. The Immigration and Nationality Act grants representation only “at no expense to the Government” \autocite{usc1362}, and the statute repeals that clause.
\item Anyone the government proposes to remove is asked, in a language they understand, whether they fear return, and told how to claim protection, before any order issues.
\item Notice of a proposed removal, or a contest of one, stays the removal until the contest is decided. The stay attaches by operation of law, not at a court's discretion, so no flight can leave while an order is being written.
\item No detained person is transferred out of reach of counsel while a contest is pending.
\item A removal carried out in breach of a stay or a court order violates the schedule, and return under the seventh term is mandatory.
\end{itemize}

Representation is the largest single difference in how these cases end. In a national study of more than a million removal cases decided from 2007 to 2012, 14 percent of detained people had a lawyer. Detained people who had one were ten and a half times as likely to succeed as those who didn't \autocite{EaglyShafer2016Counsel}. The floor's instruments require that a person lawfully in the country who faces expulsion be allowed to be represented, not that representation be paid for, so this term goes beyond them. As with the firewall, the statute can carry what the floor can't.

The third term is written against the record that opens this chapter. The Supreme Court held that people designated under the proclamation must have notice in time actually to seek habeas \autocite{jgg2025}, and that a day's notice without word of how to contest was not enough \autocite{aarp2025}. It reached those holdings after the March flights. In June 2025 it stayed a district court's order requiring written notice and a meaningful chance to raise a torture claim before removal to a third country \autocite{dvd2025}. The district court later held the policy unlawful, the First Circuit largely affirmed in September 2026, and the government has asked the Supreme Court to stay that judgment as well \autocite{dvd2026application}. A right to notice that depends on a court's remedy, when that remedy can be stayed or outrun, mostly raises costs and leaves a record. The stay by operation of law moves the waiting from the court's order into the statute.

\emph{Fourth, the firewall.} The Melt ICE Act's rule for its own grantees is extended to every federal program. Data collected to provide health care, collect taxes, educate or deliver social services may not be used to identify, locate or select anyone for immigration enforcement. This is the firewall the Council of Europe's anti-racism commission recommended a decade ago \autocite{ecri2016gpr16}. The Medicaid and IRS transfers described above are the cases it would have stopped \autocite{JoffeBlock2026MedicaidPalantir,Shaw2026IRSICE}. The firewall also covers the answer. Nothing a person submits in response to notice, or in contesting a determination, may be used to locate or select them or anyone named in it. Nor may anything held on a person's own device, or in a tool that serves them.

\emph{Fifth, the pipeline.} The contract bar is extended from detention and monitoring to the acts of the schedule, wherever a system performs them:

\begin{itemize}
\item selection on protected characteristics or their proxies;
\item use of firewalled data;
\item location by bounty;
\item scores, points tables or designations that stand in for an individual determination;
\item identification by face or other biometrics in public places;
\item screening of speech for grounds of removal;
\item identification or tracking of observers.
\end{itemize}

That reaches, as they are now written, the case-management, ImmigrationOS and ELITE contracts, the skip-tracing awards, Mobile Fortify and the speech screening described above \autocite{Koebler2025ICM,Ho2025ImmigrationOS,Cox2026ELITE,Biddle2025BountyHunters,Cox2025MobileFortify,Caputo2025CatchRevoke}. It does so without denying a successor the records it needs to make an individual determination.

The bar reaches monitoring as surveillance, not support. The Melt ICE Act's HHS grant program funds services for people affected by enforcement, delivered by organizations that have never taken part in immigration or law enforcement and that may not surveil the people they serve or pass their personal information to any federal entity \autocite{hr7190}. The statute keeps that line. Voluntary case management, legal orientation and help keeping hearing dates are support, not a monitoring program, when the person can decline them and whoever provides them may not report to enforcement.

\emph{Sixth, the record.} Preservation's no-deletion rule is applied to an agency. Before abolition takes effect, ICE's records pass to a holder adverse to the operator, under a duty to preserve them and to give access to investigators and to the people they concern. The records include case files, the pipeline's logs, scores and model versions, validation worksheets, contractor data, flight manifests, the terms of any arrangement with a foreign government to hold people removed, and body-camera and use-of-force records. Amnesty asks for prompt, independent investigation of killings by federal agents and of detention conditions it says may amount to torture \autocite{Amnesty2026Whistles}. An abolition whose last weeks resembled the NSC staff's last weeks in November 1986 would destroy the evidence those investigations need.

\emph{Seventh, remedy.} The first six terms prevent and preserve. A floor's fourth job is remedy, and the statute owes it to the people already harmed. It does so in five ways:

\begin{itemize}
\item Anyone removed in violation of the schedule or of a court order has a right of return at government expense, and their proceedings resume where they stood.
\item Families separated by removal are reunited at government expense.
\item People harmed by acts in the schedule have a cause of action for damages that sovereign immunity does not bar.
\item An independent body, outside the office that receives the agency's functions and outside the department that prosecutes, decides claims under this term. Its finding that a removal violated the schedule orders return without the successor's consent, which is what makes a review binding rather than advisory.
\item The same body investigates killings by federal agents and detention conditions during the agency's operations, with access to the records the sixth term preserves.
\end{itemize}

Amnesty asks for exactly this: prompt, effective, independent and impartial investigations into human rights violations by federal agents, and accountability for those responsible \autocite{Amnesty2026Whistles}.

Remedy that depends on the operator's cooperation is weak in the way the record shows. When the government admitted that it had removed Kilmar Abrego Garcia to El Salvador in error, the Supreme Court required it to “facilitate” his release \autocite{abregogarcia2025}. Hernández Romero's release came through a prisoner exchange, to Venezuela, and not to the proceeding he had been removed from. A right of return is the seventh term, and a finding that orders return without the operator's consent is what makes it a right.

With these terms the statute is the domestic clause's counterpart, and Table C.2 is its schedule.
